\chapter*{Global Exercise Bank --- Assertion--Reason}
\addcontentsline{toc}{chapter}{Global Exercise Bank: Assertion--Reason}
\label{chap:assertion_reason}

\begin{tcolorbox}[enhanced,colback=subtleamber,colframe=accentamber,arc=2mm,boxrule=1pt,
    title=\textbf{\large Assertion--Reason Questions (All Chapters)}]
Use the following key:\\
\textbf{(a)} Both A and R are true; R is the correct explanation of A.\\
\textbf{(b)} Both A and R are true; R is NOT the correct explanation of A.\\
\textbf{(c)} A is true; R is false.\\
\textbf{(d)} A is false; R is true.
\end{tcolorbox}

\begin{problembox}
\textbf{AR.1} \hfill \textbf{[GRB / ESS]}

\textbf{A:} The specific conductance ($\kappa$) of an electrolyte solution always decreases on dilution.

\textbf{R:} On dilution, the number of ions per unit volume decreases even though the degree of dissociation increases.
\end{problembox}
\begin{solution}
\textbf{(a)} Both A and R are true; R correctly explains A. $\kappa = \Lambda_m \cdot c/1000$: although $\Lambda_m$ increases on dilution, the concentration $c$ decreases proportionally faster, so $\kappa$ always decreases.
\end{solution}

\begin{problembox}
\textbf{AR.2} \hfill \textbf{[CNG-TH / NRJ]}

\textbf{A:} The standard hydrogen electrode (SHE) is used as the universal reference electrode.

\textbf{R:} The standard reduction potential of the SHE is arbitrarily defined as exactly zero at all temperatures.
\end{problembox}
\begin{solution}
\textbf{(c)} A is true; R is false. The SHE is the universal reference, but its potential is set to zero only at $298\,\text{K}$ as a convention. At other temperatures, the potential of the SHE is not zero (it varies with temperature due to the temperature dependence of thermodynamic quantities). The convention $E^\circ_{\text{SHE}} = 0$ applies at $298\,\text{K}$, $a_{\ce{H+}}=1$, $f_{\ce{H2}}=1\,\text{bar}$.
\end{solution}

\begin{problembox}
\textbf{AR.3} \hfill \textbf{[ATK / GRB]}

\textbf{A:} The temperature coefficient of EMF, $(\partial E/\partial T)_P$, can be positive, negative, or zero.

\textbf{R:} The sign of $(\partial E/\partial T)_P$ depends on the sign of the entropy change $\Delta S$ for the cell reaction.
\end{problembox}
\begin{solution}
\textbf{(a)} Both true; R correctly explains A. $\Delta S = nF(\partial E/\partial T)_P$. If the cell reaction proceeds with an increase in entropy ($\Delta S > 0$), then $(\partial E/\partial T)_P > 0$ (EMF increases with temperature). If $\Delta S < 0$, EMF decreases with temperature. $\Delta S = 0$ gives a temperature-independent EMF.
\end{solution}

\begin{problembox}
\textbf{AR.4} \hfill \textbf{[GRB / CNG-DPP]}

\textbf{A:} During the discharge of a lead-acid battery, the density of sulfuric acid decreases.

\textbf{R:} During discharge, sulfuric acid is consumed at both electrodes to form \ce{PbSO4}, and water is produced.
\end{problembox}
\begin{solution}
\textbf{(a)} Both true; R correctly explains A. Net discharge: $\ce{Pb + PbO2 + 2H2SO4 -> 2PbSO4 + 2H2O}$. \ce{H2SO4} (dense, $\rho=1.84\,\text{g/mL}$) is consumed; water (lighter, $\rho=1.00\,\text{g/mL}$) is produced $\Rightarrow$ density decreases.
\end{solution}

\begin{problembox}
\textbf{AR.5} \hfill \textbf{[CNG-TH / NRJ]}

\textbf{A:} \ce{KCl}, \ce{KNO3}, and \ce{NH4NO3} are preferred filling materials for salt bridges.

\textbf{R:} These salts have nearly equal ionic mobilities for cation and anion ($t_+ \approx t_-$), which minimizes liquid junction potential.
\end{problembox}
\begin{solution}
\textbf{(a)} Both true; R correctly explains A. The liquid junction potential is minimized when cation and anion transference numbers are equal. For \ce{KCl}: $t_+ = 0.491$, $t_- = 0.509$. For \ce{KNO3} and \ce{NH4NO3}: similar near-equality. These salts also have high solubility and are non-reactive with most electrolytes.
\end{solution}

\begin{problembox}
\textbf{AR.6} \hfill \textbf{[ATK / CNG-TH Advanced]}

\textbf{A:} The mobility of \ce{H+} in aqueous solution ($3.62\times10^{-7}\,\text{m}^2\,\text{V}^{-1}\,\text{s}^{-1}$) is approximately 5 times larger than a typical ion like \ce{Na+} ($5.19\times10^{-8}\,\text{m}^2\,\text{V}^{-1}\,\text{s}^{-1}$).

\textbf{R:} The anomalously high mobility of \ce{H+} is due to the Grotthuss (proton-hopping) mechanism, where \ce{H+} is transferred sequentially along a chain of hydrogen-bonded water molecules rather than migrating as a physical particle.
\end{problembox}
\begin{solution}
\textbf{(a)} Both true; R correctly explains A. In the Grotthuss mechanism, protons hop from one water molecule to the next via hydrogen-bond rearrangements (like a Newton's cradle), which is far faster than the physical diffusion of a bare ion through the viscous water medium. Similarly, \ce{OH-} has high mobility via the same mechanism running in reverse.
\end{solution}

\begin{problembox}
\textbf{AR.7} \hfill \textbf{[GRB / CNG-DPP DPP-5]}

\textbf{A:} Galvanized (zinc-coated) steel remains protected from rusting even when the zinc coating is scratched.

\textbf{R:} Tin-coated steel (tinned iron / food-grade steel) corrodes faster than uncoated iron when the tin layer is scratched.
\end{problembox}
\begin{solution}
\textbf{(b)} Both statements are true but R is not an explanation of A (they are independent facts). A is true because Zn ($E^\circ = -0.76\,\text{V}$) is more active than Fe ($-0.44\,\text{V}$) and acts as sacrificial anode. R is also true: Sn ($E^\circ = -0.14\,\text{V}$) is less active than Fe, so when the Sn coating is scratched, Fe acts as the anode and corrodes faster than bare Fe would, because it is in galvanic contact with Sn. Both are true but are separate, contrasting facts --- not a cause-effect pair.
\end{solution}
